% !TeX root = ../Book1.tex
\begin{figure}[htbp]
	\centering
	\begin{tikzcd}[row sep=large,column sep=huge]
		G \arrow[r,"{\widetilde\sigma:g\mapsto(g,f(g))}"] \arrow[d,"q"']
		& G\times W \arrow[d,"Q"] \\
		G/H \arrow[r,"{\sigma:gH\mapsto[g,f(g)]}"']
		& E_\rho
	\end{tikzcd}
	\[
		\widetilde\sigma(gh)
		=(gh,f(gh))
		=(gh,\rho(h)^{-1}f(g))
		=\widetilde\sigma(g)\cdot h.
	\]
	\caption[等变函数下降为关联丛截面]{上方映射将每个群元素及其函数值配在一起。右等变条件保证它把同一条 \(H\) 轨道送到同一条 \(H\) 轨道，因而下降为下方的截面。两个竖箭头分别是主丛商映射及关联丛商映射。}
	\label{g1:fig:associated-section-descent}
\end{figure}
